\section{Implementation}

\begin{figure}[H]
    \centering
    \includesvg[width=\columnwidth, inkscapelatex=false]{img/dqn-system-architecture.svg}
    \caption{Architecture Diagram}
    \label{fig:arch}
\end{figure}

The agent is implemented as a modular Python package (\texttt{src/dqn/}) whose components map cleanly onto the standard DQN algorithm. Figure~\ref{fig:arch} gives an overview.

\subsection{Configuration}
All hyperparameters are declared in \texttt{DQNConfig} (\texttt{dqn/config.py}), a typed dataclass whose field names match JSON config file keys one-to-one. At runtime, \texttt{DQNConfig.load} resolves a base JSON file (per environment) and applies any CLI overrides, ensuring a single source of truth for every tunable parameter.

\subsection{Policy}
Action selection is handled by the \texttt{BasePolicy} abstract class (\texttt{dqn/policies/}), from which three concrete strategies inherit:
\begin{itemize}
    \item \textbf{Epsilon-greedy} — with probability $\varepsilon$ selects a uniformly random action; otherwise acts greedily. $\varepsilon$ is annealed linearly from \texttt{explore\_start} to \texttt{explore\_end} over training.
    \item \textbf{Boltzmann} — samples from $\mathrm{softmax}(Q(s,\cdot)/T)$, where temperature $T$ is annealed analogously.
    \item \textbf{Greedy} — pure exploitation; used only during evaluation.
\end{itemize}
Each policy owns its Q-network and performs a non-blocking host-to-device copy to minimize PCIe transfer overhead during the collection phase.

\subsection{Replay Memory}
\texttt{ExperienceReplay} (\texttt{dqn/memory.py}) implements a standard ring buffer backed by pre-allocated NumPy arrays. Transitions $(s, a, r, s', d)$ are stored as contiguous blocks; when the buffer wraps around its capacity $N$ older experiences are overwritten. Sampling uses asynchronous \texttt{Tensor.copy\_} calls to overlap GPU transfer with CPU work.

\subsection{Learner}
\texttt{DQNLearner} (\texttt{dqn/learner.py}) maintains a \emph{policy network} $Q_\theta$ and a \emph{target network} $Q_{\theta^-}$ (a periodically updated copy). It implements Double DQN \cite{hasselt2016deep}: the policy network selects the greedy action in the next state, while the target network evaluates it:

\begin{equation}
    y_t = r_t + \gamma \, Q_{\theta^-}\!\left(s_{t+1},\, \arg\max_{a'} Q_\theta(s_{t+1}, a')\right)(1 - d_t)
\end{equation}

Parameters are updated via Adam with the Huber (Smooth-$L_1$) loss between $Q_\theta(s_t, a_t)$ and $y_t$. The target network is synchronized every \texttt{target\_sync\_interval} epochs; hard copies ($\tau = 1$) are used by default, with optional Polyak averaging for $\tau < 1$.

\subsection{Neural Networks}
The Q-function is approximated by a feedforward MLP (\texttt{dqn/networks/feedforward.py}) with ReLU activations. Three sizes are registered in \texttt{NetworkFactory}:

\begin{center}
\begin{tabular}{ll}
\hline
\textbf{Key} & \textbf{Hidden layers} \\
\hline
\texttt{mlp\_small}  & $128 \to 128$ \\
\texttt{mlp\_medium} & $256 \to 128 \to 64$ \\
\texttt{mlp\_large}  & $512 \to 256 \to 128 \to 64$ \\
\hline
\end{tabular}
\end{center}

Input is the 24-dimensional observation vector; output is a Q-value per discrete action.

\subsection{Training Loop}
\texttt{DQNAgent.train} (\texttt{dqn/dqn\_agent.py}) orchestrates the following phases each epoch:
\begin{enumerate}
    \item \textbf{Warmup} — random actions fill the replay buffer for \texttt{warmup\_steps} steps before any gradient updates occur.
    \item \textbf{Collection} — the current policy interacts with \texttt{num\_envs} parallel environments for \texttt{steps\_per\_epoch} steps, storing transitions in the buffer.
    \item \textbf{Optimization} — \texttt{updates\_per\_epoch} mini-batches of size \texttt{batch\_size} are sampled and used for gradient steps.
    \item \textbf{Target sync} — the target network is updated every \texttt{target\_sync\_interval} epochs.
\end{enumerate}
Exploration is annealed as a function of \texttt{current\_step / total\_steps} via the shared \texttt{linear\_anneal} utility (\texttt{dqn/logic.py}). Metrics (loss, mean episode reward, mean episode length) are logged per epoch to MLflow via \texttt{MetricsTracker}.

\subsection{Action Space Discretisation}
Since DQN requires a discrete action space and BipedalWalker-v3 natively exposes a continuous $[-1,1]^4$ action space, a wrapper (\texttt{env/bipedal\_wrapper.py}) maps a finite set of motor-speed combinations to integer indices. The discretisation scheme is described in Section~\ref{sec:bipedal-wrapper}.